\section{Limitations}
\label{sec:limitations}

\textbf{One participant.} All EEG results are for S1. The paper reports S2 and S3 at 13.59\% and
14.02\%; we did not audit them, and the dataset has two further participants. The no-EEG analysis
is cheap to repeat for them and is the first extension we would run.

\textbf{One model family and preprocessing.} We use the released model and the released
preprocessed epochs. A different architecture, frequency band or artifact treatment could find
semantic information that this pipeline does not.

\textbf{Not the authors' environment.} We ran the released configuration on Apple-silicon GPUs with
our seeds, not on the authors' hardware. Floating-point results can differ between GPU platforms; we
report the spread over seeds and the released script's own output, but cannot rule out that their
exact run reached 13.88\%.

\textbf{Two schedules.} The split comparison uses 250\,Hz and 30 epochs to fit on one machine. It is
compared only within itself, not with the 500\,Hz reproduction.

\textbf{Session holdout mixes two effects.} Category composition differs between recording days
(Appendix~\ref{app:data}), so accuracy under session holdout reflects both the absence of
within-session neighbours and a shift in label distribution. Run holdout is the cleaner control.

\textbf{Three seeds.} With $n=3$ per cell we report means, standard deviations and per-seed values,
not significance tests. Differences smaller than the seed spread should not be read as effects.

\textbf{Overlap with concurrent work.} A public code release \citep{concurrentaudit} audits
stimulus-order and run-identity confounds in Chisco's sentence-level decoding. We cite it as software;
we have not seen a paper and do not know the full extent of the overlap.
